\documentclass[8pt]{article}

\title{LegalXML}
\author{AMH}

\usepackage{outlines}
\usepackage{hyperref}
\usepackage[utf8]{inputenc}
\usepackage[ngerman]{babel}
\usepackage[T1]{fontenc}

\begin{document}
\maketitle
\tableofcontents
\newpage

\section{Zielsetzung}
Ziel dieses Projekts ist es, ein österreichisches Landesgesetz (das oberösterreichische "Landesgesetz über eine Landesabgabe für das obertägige Gewinnen mineralischer Rohstoffe
(Oö. Landschaftsabgabegesetz)")\footnote{\url{https://www.ris.bka.gv.at/GeltendeFassung.wxe?Abfrage=LrOO&Gesetzesnummer=20000943}}
zu digitalisieren. Es wird dabei eine Identität zwischen Code und Gesetz hergestellt.

\section{Generelle technische Entscheidungen}
Bei der Erstellung des Projekts werden folgende Techniken verwendet:
\begin{outline}
    \1 Für das Dependency-Management werden nix (\url{https://nixos.org/}) und
        git submodules (\url{https://git-scm.com/book/en/v2/Git-Tools-Submodules}) verwendet.
    \1 Als Programmiersprache erster Wahl wird Swift\textregistered (\url{https://www.swift.org/}) verwendet,
        sofern nötig wird (Objective) C/++ (\url{https://www.swift.org/documentation/cxx-interop/}) verwendet.
        \2 Dies deshalb, da für die Programmiersprache, in der 95\% des Codes geschrieben werden soll
            eine bekannte Programmiersprache mit Speichersicherheit verwendet werden, die gut mit C/++ interagieren kann.
            Diese Funktion ist bei Swift eingebaut (siehe Link oben).
        \2 Evaluiert wurden insofern lt. dem TIOBE-Index 2024\footnote{\url{https://www.tiobe.com/tiobe-index/}} die 20 bekanntesten Programmiersprachen:
            Python, Java, C\#, JavaScript, Go, Visual Basic, Fortran, Object Pascal/Delphi, Rust, Ruby, PHP, Swift, COBOL, Kotlin
        \2 Die offensichtlich nicht geeigneten Sprachen für Spezialanwendungen (SQL, Assembler, Scratch, MATLAB) sowie
            nicht speichersichere Sprachen (C/++) wurden ausgeschieden.
        \2 Veralterte Programmiersprache wie Fortran, COBOL, PHP, Delphi und Visual Basic werden nicht verwendet.
        \2 Nativen Programmiersprachen wurde der Vorzug gegeben, da diese eine höhere Ausführungsgeschwindigkeit erreichen
            und einfacher zu deployen sind. Java und Kotlin wären prinzipiell geeignet, jedoch erscheint die Verwendung hier unklug,
            da  die JNI-Schnittstelle zur Anbindung an C/++ wurde bereits in vergangenen Projekten an der Landesregierung verwendet
            wurde und dort für eher umständlich befunden wurde.
        \2 Damit bleiben erhalten: Go, Swift, Rust.
        \2 Lt. der offiziellen FAQ von Go ist Go schlecht geeignet um mit anderen Modulen, die in C/++ geschrieben sind,
            gelinkt zu werden (\url{https://go.dev/doc/faq#Do_Go_programs_link_with_Cpp_programs}).
        \2 Rust hat eine höhere Lernschwelle und ein komplett eigenes Ownership-Modell, das Modell der Speicherverwaltung das Swift benutzt
            (ARC \footnote{\url{https://docs.swift.org/swift-book/documentation/the-swift-programming-language/automaticreferencecounting/}}) ist
            für unsere Zwecke ausreichend schnell.
    \1 Als Build-Tool wird CMake (\url{https://cmake.org/}) verwendet.
\end{outline}

\section{Analyse des logischen Backends}
Untersucht wurde
\begin{itemize}
\end{itemize}
\endsection

\section{Module}
\begin{itemize}
    \item \textbf{ActBuilder}: Frontend für die Gestaltung von Gesetzen, basiert auf ActKit.
    \item \textbf{ActKit}: Abstrahiert den Aufbau von Gesetzen programmatisch, hält Klassen bereit für ganze Gesetze, Paragraphen, Absätze, ...
             Verwendet einen an Swift-UI \footnote{\url{https://developer.apple.com/xcode/swiftui/}} angelehnten Syntax.
    \item \textbf{PrologVM}: Zuständig für die Ausführung von Gesetzen.
\end{itemize}

Auf die Struktur der komplexeren Module wird in eigenen Abschnitten eingegangen.

\section{PrologVM}
Das eigentliche Herzstück von LegalXML ist PrologVM, dieses Modul kümmert sich um die Ausführung von Gesetzen.
Hier werden die einzelnen Paragraphen des Gesetzes auführt, und die Umsetzung erklärt.

\subsection{Gegenstand der Abgabe}

\fbox{\begin{minipage}{\textwidth}
\begin{enumerate}
    \item Das Land erhebt eine Landschaftsabgabe für das obertägige Gewinnen mineralischer Rohstoffe in Oberösterreich.
    \item Von der Erhebung ausgenommen sind:
    \begin{itemize}
        \item[-] Abraummaterial,
        \item[-] Material aus Fließgewässern, das aus flussbaulichen Gründen wieder in Fließgewässer eingebracht wird,
        \item[-] bundeseigene mineralische Rohstoffe gemäß § 4 Abs. 1 Mineralrohstoffgesetz (MinroG), BGBl. I Nr. 38/1999, in der Fassung des Bundesgesetzes BGBl. I Nr. 95/2016,
        \item[-] Kohle,
        \item[-] Material aus Seitenentnahmen und
        \item[-] Rohstoffe, deren Verwendung Maßnahmen zur Abwehr einer unmittelbaren Gefahr für das Leben oder die Gesundheit von Menschen oder zur unmittelbaren Abwehr von Katastrophen dient.
    \end{itemize}
    \item Die Gemeinde, in der sich eine Gewinnungsstätte befindet, erhält einen Ertragsanteil in Höhe von 20 \% der Landschaftsabgabe, die im Gemeindegebiet erhoben wurde. (Anm: LGBl.Nr. 96/2023)
    \item Die Landschaftsabgabe ist für Angelegenheiten des Natur- und Landschaftsschutzes sowie der Landschafts- und Ortsbildpflege, zur Verbesserung der ökologischen Infrastruktur, die Umweltbildung und Umwelterziehung sowie sonstige Maßnahmen im Bereich des Umweltschutzes zu verwenden. (Anm: LGBl.Nr. 96/2023)
    \item Die der Gemeinde gemäß Abs. 3 zufließenden Mittel sind für Angelegenheiten des Natur- und Landschaftsschutzes sowie der Landschafts- und Ortsbildpflege, zur Verbesserung der ökologischen Infrastruktur, für naturnahe Erholungsformen in der Gemeinde, die Umweltbildung oder die Umwelterziehung zu verwenden. (Anm: LGBl.Nr. 96/2023)
    \item Die Überweisung des Ertragsanteils an die Gemeinden hat jeweils spätestens am 30. Juni für das vorangegangene Kalenderjahr zu erfolgen. (Anm: LGBl.Nr. 96/2023)
\end{enumerate}
\end{minipage}}

Im ersten Paragraph definiert eine neue Abgabe, die dann in den folgenden Pragraphen und Absätzen genauer ausgestaltet wird.

\subsection{Begriffsbestimmungen}

\fbox{\begin{minipage}{\textwidth}
Im Sinn des Landesgesetzes ist:
\begin{enumerate}
    \item „Abraummaterial“: jedes beim Gewinnen anfallende, nicht verwertbare Material (zB taubes Gestein, Abschlämmbares) sowie Materialien,
    die zur Böschungsherstellung, Rekultivierung oder Geländegestaltung (zB Lärmschutz- oder Hochwasserschutzdämme) betriebsintern verwendet werden;
    \item „Betreiberin“ bzw. „Betreiber“: jede physische und juristische Person sowie jeder sonstige Rechtsträger, die bzw. der ein Gewinnen gewerblich oder berufsmäßig durchführt;
    \item „Gewinnen“: das Lösen oder Freisetzen (Abbau) mineralischer Rohstoffe einschließlich der durch dieselbe Betreiberin bzw. denselben Betreiber
            vorgenommenen damit zusammenhängenden vorbereitenden, begleitenden und nachfolgenden Tätigkeiten zur Aufbereitung des Naturmaterials;
    \item „Gewinnungsstätte“: Steinbruch bzw. Entnahmestelle von mineralischen Rohstoffen unabhängig davon, ob dafür eine Bewilligungspflicht nach dem MinroG besteht;
    \item „Mineralischer Rohstoff“: jedes Mineral, Mineralgemenge oder Gestein (Fest- und Lockergestein), wenn es natürlicher Herkunft ist;
    \item „Seitenentnahme“: obertägiges Gewinnen im direkten Areal eines Bauprojekts zwecks Verwendung bei diesem Bauprojekt;
    \item „Verwertung“: Übergabe an Dritte oder betriebsinterne Übergabe zur Weiterverarbeitung nach der Aufbereitung.
\end{enumerate}
\end{minipage}}

\subsection{Abgabepflichtige bzw. Abgabepflichtiger}

\fbox{\begin{minipage}{\textwidth}
Abgabepflichtige bzw. Abgabepflichtiger ist die Betreiberin bzw. der Betreiber einer Gewinnungsstätte eines abgabepflichtigen Materials.
\end{minipage}}

\subsection{Abgabenbefreiung}

\fbox{\begin{minipage}{\textwidth}
Von der Landschaftsabgabe befreit sind Betreiberinnen bzw. Betreiber, deren Abgabenschuld im jeweiligen Kalenderjahr weniger als 120 Euro beträgt.
\end{minipage}}

\subsection{Höhe der Abgabe}

\fbox{\begin{minipage}{\textwidth}
\begin{enumerate}
    \item Die Höhe der Landschaftsabgabe beträgt 15,95 Cent pro Tonne gewonnenen und verwerteten mineralischen Rohstoffs.
    \item Der im Abs. 1 festgesetzte Tarif ändert sich jeweils zum 1. Jänner entsprechend den durchschnittlichen Änderungen
            des von der Bundesanstalt „Statistik Austria“ für das zweitvorangegangene Jahr verlautbarten Verbraucherpreisindex 2015 oder eines an seine Stelle
            tretenden Index. Bezugsgröße für die erstmalige Änderung zum Stichtag 1. Jänner 2025 ist der durchschnittliche Indexwert für das Jahr 2017;
            Bezugsgröße für jede weitere Änderung ist der durchschnittliche Indexwert des der jeweils letzten Änderung zweitvorangegangenen Kalenderjahres.
            Ein sich aus dieser Berechnung ergebender neuer Betrag ist auf einen vollen Hundertstel-Centbetrag zu runden, wobei Beträge bis einschließlich
            0,005 Cent abgerundet und Beträge über 0,005 Cent aufgerundet werden. Eine solchermaßen ermittelte Änderung des Tarifs wird nur dann wirksam,
            wenn der geänderte Betrag von der Landesregierung vor dem Stichtag 1. Jänner im Landesgesetzblatt für Oberösterreich kundgemacht wurde. (Anm: LGBl.Nr. 95/2022, 96/2023)
\end{enumerate}
\end{minipage}}

\subsection{Entstehen der Abgabenschuld}

\fbox{\begin{minipage}{\textwidth}
Die Abgabenschuld entsteht zu dem Zeitpunkt, in dem das gewonnene Material verwertet wird.
\end{minipage}}

\subsection{Aufzeichnungspflicht}

\fbox{\begin{minipage}{\textwidth}
Die bzw. der Abgabepflichtige ist verpflichtet, zur Feststellung der Abgabe und der Grundlagen ihrer Berechnung Aufzeichnungen zu führen.
\end{minipage}}

\subsection{Anzeigepflicht}

\fbox{\begin{minipage}{\textwidth}
Die bzw. der Abgabepflichtige hat den Beginn und das Ende eines abgabepflichtigen Gewinnens binnen vier Wochen der Abgabenbehörde anzuzeigen.
\end{minipage}}

\subsection{Selbstbemessung, Fälligkeit}

\fbox{\begin{minipage}{\textwidth}
\begin{enumerate}
    \item Die bzw. der Abgabepflichtige hat die Abgabe selbst zu bemessen. Die Abgabenerklärung ist jeweils
            bis 30. April eines jeden Jahres (Fälligkeitstag) für die im Vorjahr entstandene Abgabenschuld einzureichen.
    \item Die Abgabenerklärung ist nach Gemeinden und nach Gewinnungsstätten aufzugliedern.
            Die bzw. der Abgabepflichtige hat den Abgabenbetrag zu berechnen und die Abgabe am Fälligkeitstag zu entrichten. (Anm: LGBl.Nr. 96/2023)
\end{enumerate}
\end{minipage}}

\subsection{Behörde}

\fbox{\begin{minipage}{\textwidth}
Abgabenbehörde ist die Landesregierung.
\end{minipage}}

\subsection{Übermittlung von Daten}

\fbox{\begin{minipage}{\textwidth}
\begin{enumerate}
    \item Die Abgabenbehörde hat jährlich jeweils bis 30. September eines jeden Jahres die ihr im Rahmen der Vorlage von Abgabenerklärungen für die im Vorjahr entstandene Abgabenschuld bekannt gegebenen
    \begin{itemize}
        \item[-] Namen oder Firmenbezeichnungen oder Firmenbuchnummern der Betreiberinnen und Betreiber von Gewinnungsstätten,
        \item[-] den Ort der Gewinnung (politischer Bezirk und Gemeinde) und
        \item[-] die Menge des dort gewonnenen und verwerteten mineralischen Rohstoffs
    \end{itemize}
    in elektronischer Form an die für die Vollziehung des Oö. Raumordnungsgesetzes 1994 (Oö. ROG 1994) zuständige Behörde und
    die die Parteistellung gemäß § 81 MinroG wahrnehmende Dienststelle des Amtes der Landesregierung zum Zweck der Einspeisung in das Oö. Rohstoffinformationssystem zu übermitteln.
    \item Sollte ein Hinderungsgrund für die vollständige oder auch nur teilweise Datenübermittlung vorliegen, hat die Abgabenbehörde die für die Vollziehung des Oö. ROG 1994
            zuständige Behörde und die die Parteistellung gemäß § 81 MinroG wahrnehmende Dienststelle des Amtes der Landesregierung hierüber unter Mitteilung des Hinderungsgrundes zu informieren.
    \item Die Naturschutzbehörden haben in Bewilligungsbescheide gemäß § 5 Z 11 und 15 Oö. Natur- und Landschaftsschutzgesetz 2001 einen Hinweis
        auf die Verpflichtungen nach dem Oö. Landschaftsabgabegesetz aufzunehmen. (Anm: LGBl.Nr. 96/2023)
\end{enumerate}
\end{minipage}}

\subsection{Schlussbestimmungen}

\fbox{\begin{minipage}{\textwidth}
    \begin{enumerate}
        \item Dieses Landesgesetz tritt mit dem seiner Kundmachung im Landesgesetzblatt für Oberösterreich folgenden Monatsersten in Kraft.
        \item Personen, die zum Zeitpunkt des Inkrafttretens dieses Landesgesetzes bereits eine Gewinnungsstätte eines abgabepflichtigen Materials betreiben,
                haben ihre Tätigkeit binnen vier Wochen nach Inkrafttreten dieses Landesgesetzes der Abgabenbehörde anzuzeigen.
        \item Abweichend von § 9 ist die Abgabenerklärung für das erste Halbjahr 2018 bis längstens 31. Oktober 2018 einzureichen und auch
                die Abgabe für das erste Halbjahr 2018 bis längstens 31. Oktober 2018 zu entrichten.
    \end{enumerate}
\end{minipage}}



\end{document}